% Input after \appendix. Requires amsmath, amssymb, and natbib.
% Bibliography entries are in paper/theory_sources.bib.
\section{Identification, restricted decision risk, and measurement assumptions}
\label{app:theory}

This appendix gives an exact counterexample and the assumptions under which
additional measurements resolve it. The example is a construction, not an
estimated model of human preferences. Its decision objective is one explicit
choice among objectives for influenceable preferences
\citep{carroll2024}. Neither the counterexample nor the sampling
identity below establishes a new general estimation method, human prevalence,
or a neural learning result.

\subsection{Two mechanisms with a common observation law}
\label{app:theory:witness}

Let a user's initial binary state be $Z_0\sim\operatorname{Bernoulli}(p)$.
A logging policy draws $A\in\{0,1\}$ with
$\Pr(A=1\mid Z_0=z)=e_z$. The post-action state is
$(P_1,S_1)$, where $P_1$ is the persistent component and $S_1$ is an expressed
stance. The initial state is $(Z_0,Z_0)$, and the common, action-independent
observation function is $O=S_1$. Fix copying probabilities $c_0,c_1\in[0,1]$.
If $A=Z_0$, both mechanisms leave the state unchanged. If $A\ne Z_0$,
the expression mechanism $E$ and the transition mechanism $T$ use
\begin{align}
 E:\quad (P_1,S_1)&=
 \begin{cases}(Z_0,A),&\text{with probability }c_A,\\
 (Z_0,Z_0),&\text{with probability }1-c_A;\end{cases}\\
 T:\quad (P_1,S_1)&=
 \begin{cases}(A,A),&\text{with probability }c_A,\\
 (Z_0,Z_0),&\text{with probability }1-c_A.\end{cases}
\end{align}
The labels ``persistent'' and ``expression'' are part of the construction.
They do not turn a recorded utterance into an independently measured state.

\paragraph{Proposition 1 (immediate observational equivalence).}
For any $p,e_0,e_1,c_0,c_1$, the joint law of $(Z_0,A,O)$ is identical under
$E$ and $T$. This remains true even if the historical log records $Z_0$ without
error and the action has randomized overlap.

\paragraph{Proof.}
When $a=z$, $\Pr(O=z\mid Z_0=z,A=a)=1$ in both mechanisms. When $a\ne z$,
both give $\Pr(O=a\mid Z_0=z,A=a)=c_a$ and
$\Pr(O=z\mid Z_0=z,A=a)=1-c_a$. Multiplying by the same initial-state and
logging probabilities proves equality of every joint atom. Hence the total
variation distance is zero. Any independent sample of historical logs, and
any statistic or randomized post-processing of that sample, has the same
law in the two worlds. This is equality by construction, not a test passed
by human data.\hfill$\square$

For a fresh user drawn independently from the same initial-state population,
define the value of an intervention that always chooses action $a$ by
\begin{equation}
 V_m(a)=\Pr_m\{a=P_1(a)\},\qquad m\in\{E,T\}.
 \label{eq:theory:persistent-value}
\end{equation}
The evaluated action itself may change the state against which it is scored.
This objective rewards post-action matching; it is not a definition of welfare
or a requirement to preserve the initial state. Direct evaluation gives
\begin{align}
 (V_E(0),V_E(1))&=(1-p,p),\\
 (V_T(0),V_T(1))&=(1-p+pc_0,\ p+(1-p)c_1).
\end{align}
For the exact witness $p=3/5$, $c_0=1$, $c_1=0$, these values are
\begin{equation}
 (V_E(0),V_E(1))=(2/5,3/5),\qquad
 (V_T(0),V_T(1))=(1,3/5).
 \label{eq:theory:witness-values}
\end{equation}
Thus action $1$ is uniquely optimal among constant actions in $E$, while
action $0$ is uniquely optimal in $T$. With $e_0=e_1=1/2$, the nonzero
immediate-log atoms $(z,a,o)$ have masses
\begin{equation}
 (0,0,0):1/5,\quad (0,1,0):1/5,\quad
 (1,0,0):3/10,\quad (1,1,1):3/10.
\end{equation}
The same equivalence also holds for the separately archived logging design
$(e_0,e_1)=(7/20,7/10)$; the logging designs must not be interchanged when
reporting a delayed-log total variation distance.

\subsection{The corrected expected-risk minimax statement}
\label{app:theory:minimax}

A learner receives $n$ independent historical immediate logs $D_n$ and may
use independent randomization $U$. It outputs one action $\delta(D_n,U)$ for
a fresh anonymous user. The fresh user's $Z_0$ is unavailable at deployment
and independent of $D_n,U$. The comparison class contains the two constant
actions. Define expected regret, averaging over training data and learner
randomization, by
\begin{equation}
 R_m(\delta)=\max_{a\in\{0,1\}}V_m(a)
             -\mathbb E_m[V_m(\delta(D_n,U))].
\end{equation}
This restriction is essential: it is not a lower bound for all personalized
policies with access to the current user's state.

\paragraph{Proposition 2 (two-point minimax risk).}
For the witness in Eq.~\eqref{eq:theory:witness-values}, for every $n\geq0$,
\begin{equation}
 \inf_\delta\max\{R_E(\delta),R_T(\delta)\}=\frac{2}{15},
 \label{eq:theory:minimax}
\end{equation}
where the infimum allows randomized learners. It is attained by ignoring
the data and choosing action $1$ with probability $1/3$.

\paragraph{Proof.}
By Proposition 1, $q=\Pr_m\{\delta(D_n,U)=1\}$ is the same under both
mechanisms. Equation~\eqref{eq:theory:witness-values} therefore gives
\begin{equation}
 R_E(\delta)=\frac15(1-q),\qquad
 R_T(\delta)=\frac25q.
\end{equation}
If $q\leq1/3$, the first expression is at least $2/15$; if $q\geq1/3$,
the second is at least $2/15$. At $q=1/3$ they are equal. The data-independent
randomized rule attains this probability for every sample size.\hfill$\square$

\paragraph{Correction: deterministic data dependence is not a fixed action.}
The minimum worst-case regret over the two \emph{fixed} deterministic actions
is $1/5$. The stronger assertion that every deterministic learning algorithm
has expected worst-case regret at least $1/5$ is false. For $n\geq1$, consider
the deterministic rule
\begin{equation}
 \delta(D_n)=\mathbf 1\{Z_{0,1}=0\},
\end{equation}
where $Z_{0,1}$ is the first historical user's initial state. In either world,
$q=2/5$. Its two expected regrets are $3/25$ and $4/25$, with maximum
$4/25<1/5$. The random historical sample supplies randomness even though the
algorithm has no random seed. Proposition 2 remains a lower bound for this
algorithm; no claim that the deterministic subclass attains $2/15$ at each
fixed finite $n$ is needed.

\paragraph{Observed-current-state escape.}
If the fresh user's $Z_0$ is observed before acting, the policy $A=Z_0$
causes no copying branch and has $A=P_1=Z_0$ with probability one under both
mechanisms. Its value is one and its regret against the optimal policy in
that enlarged class is zero. Historical mechanism nonidentification survives,
but the restricted anonymous-user decision lower bound no longer applies.
Likewise, the witness alone does not establish impossibility for an unknown
longitudinal system with other informative interventions or observations.

\subsection{What an additional measurement must identify}
\label{app:theory:measurement}

In the binary construction, define an ideal neutral probe by the state map
$(P_1,S_1)\mapsto(P_1,P_1)$ followed by the same emission function. Its
output is $B=P_1$. The only immediate-log atom whose delayed outcome differs
between the witness worlds is $(Z_0,A,O)=(1,0,0)$: it has $B=1$ in $E$ and
$B=0$ in $T$. Consequently, the total variation distance between the extended
laws of $(Z_0,A,O,B)$ is $3/10$ under balanced randomized logging, and
$(3/5)(1-7/10)=9/50$ under the archived state-dependent logging design.
The probe identifies these two stipulated worlds because its exclusion and
measurement properties were specified, not because time elapsed.

The following finite-state statement makes the measurement requirement
explicit. Let $Z_0$ and $P_1$ take values in $\{1,\ldots,r\}$ and let
$T_a(z,j)=\Pr(P_1=j\mid Z_0=z,\operatorname{do}(A=a))$. Assume that:
\begin{enumerate}
 \item the baseline states being conditioned on are observed without error,
 have positive probability, and the logged action is randomized conditional
 on these states with positive probability for each relevant action;
 \item the probe does not change $P_1$, which remains stable over the specified
 horizon; and its known measurement law is the row-stochastic matrix
 $M_{jb}=\Pr(B=b\mid P_1=j)$, with no remaining direct dependence on action,
 baseline, or transient expression conditional on $P_1$;
 \item the full-population conditional probe distributions are observed or
 recoverable under an explicitly justified sampling and response model.
\end{enumerate}
Randomizing invitations to a delayed survey does not by itself justify the
third condition when actual response depends on the unobserved outcome.

\paragraph{Proposition 3 (known-emission identification).}
Writing $Q_a(z,b)=\Pr(B=b\mid Z_0=z,A=a)$, these assumptions give
$Q_a=T_aM$. If $M$ has full row rank $r$, then
\begin{equation}
 T_a=Q_aM^+,\qquad
 \|\widehat T_a-T_a\|_F
 \leq\|\widehat Q_a-Q_a\|_F\,\|M^+\|_2,
 \quad \widehat T_a=\widehat Q_aM^+,
 \label{eq:theory:known-emission}
\end{equation}
where $M^+$ is its Moore--Penrose inverse. The error bound assumes $M$ is
known exactly; it does not account for estimating or misspecifying $M$.

\paragraph{Proof.}
Conditioning on $P_1$ gives $Q_a=T_aM$. Full row rank implies $MM^+=I_r$,
so right multiplication identifies $T_a$. Submultiplicativity of the
Frobenius and spectral norms proves the error bound.\hfill$\square$

Small singular values can make finite-sample inversion unstable. Full row
rank is a sufficient condition here, not a claim that every target functional
requires full transition identification. Conversely, for
\begin{equation}
 M=\begin{pmatrix}1&0\\0&1\\0&1\end{pmatrix},\qquad
 T^{(1)}=I_3,\qquad
 T^{(2)}=\begin{pmatrix}1&0&0\\0&0&1\\0&1&0\end{pmatrix},
\end{equation}
$T^{(1)}M=T^{(2)}M$ although $T^{(1)}\ne T^{(2)}$. Thus additional probe
samples alone do not fix a measurement channel that aliases these states.
An unknown emission matrix also leaves a factorization problem; this theorem
does not learn or validate it from arbitrary conversational logs.

\subsection{The standard residual estimator and its limits}
\label{app:theory:residual}

For one randomized action stratum, let the finite population of $N$ logged
units have immediate agreement scores $O_i\in[0,1]$ and delayed scores
$B_i\in[0,1]$. The target here is the finite-population mean
$\overline B_N=N^{-1}\sum_i B_i$. Let $S$ be a uniformly selected subset of
$k$ units, sampled without replacement, on which both scores are observed.
Then the difference estimator is
\begin{equation}
 \widehat\mu_{\rm diff}=\overline O_N+
 \frac1k\sum_{i\in S}(B_i-O_i).
 \label{eq:theory:difference}
\end{equation}
It is standard model-assisted estimation, closely related to the mean
estimator in prediction-powered inference
\citep{ppi,mozer2026difference}.

\paragraph{Proposition 4 (design expectation).}
Conditional on the complete finite population,
$\mathbb E_S[\widehat\mu_{\rm diff}]=\overline B_N$. More generally, if
$R_i$ indicates that unit $i$ is measured and its positive inclusion
probability $\pi_i=\Pr(R_i=1\mid\{O_j,B_j\}_{j=1}^N)$ is known, then
\begin{equation}
 \widehat\mu_{\rm HT-diff}=\frac1N\sum_{i=1}^N
 \left[O_i+\frac{R_i}{\pi_i}(B_i-O_i)\right]
 \label{eq:theory:ht-difference}
\end{equation}
has the same conditional expectation. Dependence among the $R_i$ does not
change this expectation identity, although it matters for variance.

\paragraph{Proof.}
In the first design, $\mathbb E[\mathbf1\{i\in S\}]=k/N$; linearity gives
$\mathbb E[k^{-1}\sum_{i\in S}(B_i-O_i)]
=N^{-1}\sum_i(B_i-O_i)$. In the second design,
$\mathbb E[R_i/\pi_i]=1$ for each unit. Adding $\overline O_N$ yields
$\overline B_N$ in each case.\hfill$\square$

For a superpopulation response model using observed covariates $W$, a
corresponding sufficient condition is
$\mathbb E[R\mid B,O,W]=\pi(W)>0$ with known $\pi(W)$. A randomized
invitation probability is not the full measurement-inclusion probability
if nonresponse is selective. Plug-in estimated weights, self-normalization,
outcome-dependent selection with unknown probabilities, and fitting a proxy
on the same sampled labels require their own analysis; Proposition 4 does
not supply a blanket finite-sample unbiasedness guarantee for them.

\paragraph{Clipping changes the expectation.}
The archived first finite-state implementation clips
$\widehat\mu_{\rm diff}$ to $[0,1]$. This can reduce squared error for a
target in $[0,1]$, but does not preserve unbiasedness. For the finite
population $N=2$, $k=1$, $(O_1,O_2)=(1,0)$ and $(B_1,B_2)=(0,0)$,
Eq.~\eqref{eq:theory:difference} equals $-1/2$ or $1/2$ with equal
probability. Its expectation is the true mean zero. After clipping, the
expectation is $1/4$. This counterexample concerns estimator expectation;
it does not invalidate a properly reported Monte Carlo regret comparison.

\paragraph{Unbiasedness does not imply a decision advantage.}
For independent identically distributed pairs with finite second moments,
the same $k$ uniformly selected labels from $N$ pairs, and coefficient one,
direct covariance expansion gives
\begin{equation}
 \operatorname{Var}(\widehat\mu_{\rm diff})
 =\frac{\operatorname{Var}(B)}{k}
 +\left(\frac1k-\frac1N\right)
 \left[\operatorname{Var}(O)-2\operatorname{Cov}(B,O)\right].
\end{equation}
Here $\operatorname{Var}(\overline O_N)=\operatorname{Var}(O)/N$,
$\operatorname{Var}(\overline{B-O}_S)=\operatorname{Var}(B-O)/k$, and
$\operatorname{Cov}(\overline O_N,\overline{B-O}_S)
=\operatorname{Cov}(O,B-O)/N$, which proves the formula. Thus the estimator
need not beat the ordinary mean of the same $k$ delayed labels. Smaller
estimation error also need not change which action is selected, and a biased
estimator can have smaller decision regret on a particular population.
All same-label anchor means and simple shrinkage/mixture baselines therefore
belong in a method comparison. In the constructed truthful-transition world
$B_i=O_i$ for every pair, so its correction is identically zero; observing
this equality in a replay is an algebra check, not an independent empirical
discovery about human behavior.

\subsection{The empirical target is a separate estimand}
\label{app:theory:bridge}

The constructed EndoPAHF assay assigns a historical interaction an immediate
label and a stipulated delayed label. It then scores a later model answer
against that fixed label. No state-transition kernel is applied to the
\emph{evaluated} later answer. Its expected correct-label score has the form
\begin{equation}
 U_B(\theta)=\mathbb E_{(X,B)\sim\mathcal D}
                   [\pi_\theta(B\mid X)],
 \label{eq:theory:predictive-target}
\end{equation}
with a fixed constructed distribution $\mathcal D$. If $\pi_\theta$ is a
full-vocabulary next-token distribution, all tokens other than the native
correct answer token receive reward zero; a constrained option-normalized
policy instead defines a different decision rule and must be labeled as such.
Neither version equals Eq.~\eqref{eq:theory:persistent-value} merely because
the label is called delayed. Established users appearing in both learning
and evaluation also do not become unseen users when the query prefix changes.

Replacing $(B_i,O_i)$ in the residual identity by fixed, differentiable
delayed and immediate losses can identify the corresponding sampled loss
or its gradient when the same sampling conditions hold and differentiation
commutes with expectation. It does not prove that this gradient improves
Eq.~\eqref{eq:theory:predictive-target}, let alone an action-dependent human
state objective. Feedback generated from the learner's own policy adds
distributional derivatives not supplied by a fixed-target identity.
The distinction between likelihood, held-out feedback gradients, and value
is already central to \citet{privilegedlikelihood}.

Finally, an immediate-versus-delayed human survey study can identify effects
on a specified measured response under randomization and justified attrition
assumptions. It cannot establish $B=P_1$, a known full-rank $M$, absence of
measurement reactivity, or a uniquely correct welfare objective simply by
having a follow-up date. Those are separate validation and normative tasks.
